\chapter{Factor graphs and belief propagation}
\label{ch:factor-graphs-bp}
% colours of the paper's house palette (tools/figstyle.py), defined locally so this part
% compiles inside the compendium preamble, which does not load the paper's main.tex
\definecolor{bpBlue}{HTML}{2A78D6}
\definecolor{bpInk}{HTML}{52514E}
\definecolor{bpFill}{HTML}{E1E0DC}

The paper's central computational claim is that the score of a non-Gaussian,
correlated law in $\nsites$ dimensions is available exactly, by ``one forward--backward
sweep of functional message passing''. This chapter builds the objects in that sentence
from scratch: what a factor graph is, what the sum-product algorithm does on it, why the
algorithm is exact when the graph is a tree, and why conditioning the chain prior on
coordinatewise noised observations leaves a tree. It then specialises the general
algorithm to the two message sweeps of the paper's equation for $\lmsg_i$ and $\rmsg_i$
(Appendix B, the equation labelled \emph{bp}), reads off the site marginals, the pairwise
marginals and the evidence, and places the recursion next to its two classical
relatives, the forward--backward algorithm for hidden Markov models and the Kalman
smoother.

\section{Graphical models and factor graphs}

\subsection{Factorised functions}

Let $a=(a_1,\dots,a_\nsites)$ be a vector of variables, each taking values in a set
$\mathcal{A}$ (a finite alphabet, or $\Reals$). A \term{factorisation} of a
non-negative function $f$ of $a$ is a representation
\begin{equation}
  f(a)=\prod_{\alpha\in\mathcal{F}}\psi_\alpha(a_{\partial\alpha}),
  \label{eq:bp-factorisation}
\end{equation}
where $\mathcal{F}$ is a finite index set of \term{factors}, each factor $\alpha$ has a
\term{scope} $\partial\alpha\subseteq\{1,\dots,\nsites\}$, and $a_{\partial\alpha}$ is the
sub-vector of the variables in that scope. The \term{factor graph} of the factorisation
is the bipartite graph with one \term{variable node} per index $i$, one \term{factor
node} per factor $\alpha$, and an edge between $i$ and $\alpha$ whenever
$i\in\partial\alpha$. A factor with a scope of size one is called
\term{unary}; a scope of size two is \term{pairwise}.

When $f$ is a probability density up to a constant, $p(a)=f(a)/Z$ with
$Z=\int f(a)\dd a$, the factor graph is a \term{graphical model} of $p$. The constant
$Z$ is the \term{normaliser}, also called the \term{partition function} in physics and,
when $f$ is a joint density of data and latent variables evaluated at the observed data,
the \term{evidence} or \term{marginal likelihood}.

\begin{example}[The chain prior]
\label{ex:bp-chain-prior}
The paper's clean law, its Section 2, equation \emph{prior}, is
$\Pz(\va)=\mu(a_1)\prod_{i=2}^{\nsites}\kernel(a_i\mid a_{i-1})$. It is a factorisation
with one unary factor $\mu$ of scope $\{1\}$ and $\nsites-1$ pairwise factors
$\kernel(\cdot\mid\cdot)$ of scopes $\{i-1,i\}$. Its factor graph is a path: variable
nodes $a_1,a_2,\dots,a_\nsites$ with a factor node between each consecutive pair, and
a dangling unary factor at $a_1$. Here $Z=1$ because each $\kernel(\cdot\mid a)$
integrates to one.
\end{example}

\subsection{Conditional independence}

The factor graph encodes which variables are conditionally independent of which. The
statement we need is the following. If removing the variable node $a_i$ (together with
its edges) disconnects the graph into components $\mathcal{C}_1,\dots,\mathcal{C}_r$,
then under $p$ the groups of variables in different components are conditionally
independent given $a_i$. The proof is one line: every factor lies inside one component
(its scope, minus $i$, is connected through the factor node), so
$f(a)=\prod_j g_j(a_{\mathcal{C}_j},a_i)$, and a joint density that is a product of
functions of the separate groups (each also depending on $a_i$) is, given $a_i$, a
product density.

Applied to the path of \cref{ex:bp-chain-prior}, removing $a_i$ separates the past
$\{a_1,\dots,a_{i-1}\}$ from the future $\{a_{i+1},\dots,a_\nsites\}$: this is the
\term{Markov property}, and it is the only use we make of the word ``Markov''.

\subsection{Trees, and why marginalisation is hard in general}

A factor graph is a \term{tree} when it is connected and has no cycles. A path is a
tree. The \term{marginal} of $p$ at a variable $i$ is
$p_i(a_i)=\frac1Z\int f(a)\dd a_{\setminus i}$, integrating over every variable but the
$i$-th. Written out on a finite alphabet of size $\ngrid$, the integral is a sum over
$\ngrid^{\nsites-1}$ configurations, which for the paper's $\ngrid=401$ and $\nsites=32$
is not a computation. Every algorithm in this chapter is a way of organising that sum
so that it costs $O(\nsites\ngrid^2)$ instead, and the organisation is possible exactly
because the graph is a tree.

\section{The sum-product algorithm}

\subsection{Messages}

The \term{sum-product algorithm}, also called \term{belief propagation}, passes two
kinds of functions along the edges of the factor graph. For an edge between variable
$i$ and factor $\alpha$:
\begin{align}
  \text{variable to factor:}\qquad
  m_{i\to\alpha}(a_i)&=\prod_{\beta\in\partial i\setminus\{\alpha\}} m_{\beta\to i}(a_i),
  \label{eq:bp-v2f}\\
  \text{factor to variable:}\qquad
  m_{\alpha\to i}(a_i)&=\int\psi_\alpha(a_{\partial\alpha})
    \prod_{j\in\partial\alpha\setminus\{i\}} m_{j\to\alpha}(a_j)\dd a_{\partial\alpha\setminus i},
  \label{eq:bp-f2v}
\end{align}
where $\partial i$ is the set of factors whose scope contains $i$. An empty product is
$1$, so a leaf variable (one attached to a single factor) sends the constant message
$1$, and a leaf factor (a unary factor) sends its own function
$m_{\alpha\to i}(a_i)=\psi_\alpha(a_i)$. The names are literal: a message is a function
of the variable at the edge's variable end, and \cref{eq:bp-f2v} sums (integrates) a
product.

On a tree, every message can be computed once and in a definite order: root the tree
anywhere, compute messages from the leaves towards the root, then from the root back
towards the leaves. Each edge carries one message in each direction, and each message
is computed after all the messages it depends on. The \term{belief} at a variable and at
a factor are then
\begin{equation}
  b_i(a_i)\propto\prod_{\alpha\in\partial i} m_{\alpha\to i}(a_i),\qquad
  b_\alpha(a_{\partial\alpha})\propto\psi_\alpha(a_{\partial\alpha})
   \prod_{j\in\partial\alpha} m_{j\to\alpha}(a_j).
  \label{eq:bp-beliefs}
\end{equation}

\subsection{Exactness on a tree}

\begin{theorem}[Sum-product is exact on a tree]
\label{thm:bp-exact}
Let the factor graph of \cref{eq:bp-factorisation} be a tree. Then the beliefs of
\cref{eq:bp-beliefs} are the exact marginals: $b_i=p_i$ and $b_\alpha=p_{\partial\alpha}$,
where $p_{\partial\alpha}$ is the marginal of $p$ on the scope of $\alpha$. Moreover, for
every variable $i$, $Z=\int\prod_{\alpha\in\partial i} m_{\alpha\to i}(a_i)\dd a_i$.
\end{theorem}

\begin{proof}
The proof is an induction on subtrees, and it rests on one fact about integrals of
products: the \term{distributive law}, $\int g(u)h(v)\dd u\dd v=\int g\dd u\int h\dd v$
when the two factors depend on disjoint sets of variables.

Fix an edge $(i,\alpha)$ and delete it. Because the graph is a tree, this splits it into
two subtrees; call $T_{\alpha\to i}$ the one containing $\alpha$ and $T_{i\to\alpha}$ the one
containing $i$. Write $F_{\alpha\to i}(a_i)$ for the integral, over every variable of
$T_{\alpha\to i}$ other than $a_i$, of the product of all factors in $T_{\alpha\to i}$; and
$F_{i\to\alpha}(a_i)$ likewise for $T_{i\to\alpha}$ (the factors on the $i$ side other than
$\alpha$). We claim
\begin{equation}
  m_{\alpha\to i}=F_{\alpha\to i},\qquad m_{i\to\alpha}=F_{i\to\alpha}.
  \label{eq:bp-claim}
\end{equation}
The claim is proved by induction on the number of nodes in the subtree behind the
message. If $T_{\alpha\to i}$ consists of $\alpha$ alone, then $\alpha$ is unary and both
sides equal $\psi_\alpha(a_i)$. If $T_{i\to\alpha}$ consists of $i$ alone, both sides are
$1$. Otherwise, suppose the claim holds for all smaller subtrees. The subtree
$T_{\alpha\to i}$ is $\alpha$ together with, for each $j\in\partial\alpha\setminus\{i\}$, the
subtree $T_{j\to\alpha}$; these subtrees are disjoint (no cycles) and share no variables
with one another (a shared variable would be a cycle through $\alpha$). Hence the product
of all factors in $T_{\alpha\to i}$ is $\psi_\alpha\prod_j(\text{factors of }T_{j\to\alpha})$,
and integrating out the variables of each $T_{j\to\alpha}$ other than $a_j$ can be done
separately by the distributive law, giving $\psi_\alpha\prod_j F_{j\to\alpha}(a_j)$; the
remaining integral over $a_{\partial\alpha\setminus i}$ is exactly \cref{eq:bp-f2v} once
$F_{j\to\alpha}=m_{j\to\alpha}$ by the induction hypothesis. The variable-to-factor case
is the same argument with the subtrees $T_{\beta\to i}$, $\beta\in\partial i\setminus\{\alpha\}$,
which share only the variable $a_i$; no integral is needed and
\cref{eq:bp-v2f} follows.

Now take a variable $i$. The whole graph is the union of the subtrees $T_{\alpha\to i}$,
$\alpha\in\partial i$, which pairwise share only $a_i$. By the distributive law again,
$\int f(a)\dd a_{\setminus i}=\prod_{\alpha\in\partial i}F_{\alpha\to i}(a_i)$, and by
\cref{eq:bp-claim} this is the unnormalised belief. Integrating once more over $a_i$
gives $Z$. For a factor $\alpha$, the whole graph is $\alpha$ together with the subtrees
$T_{j\to\alpha}$, $j\in\partial\alpha$, sharing no variables, and the same reasoning gives
$b_\alpha=p_{\partial\alpha}$.
\end{proof}

Three remarks on the theorem. First, nothing in the proof depends on the variables
being discrete: on $\Reals$ the messages are functions and the sums are integrals,
which is what the paper means by \emph{functional} message passing. Second, the
proof shows that the belief at a factor integrates, over the variables of the factor
other than $i$, to the belief at $i$: beliefs are \term{locally consistent}. Third,
on a graph with cycles the messages of \cref{eq:bp-v2f,eq:bp-f2v} still make sense
and can be iterated (\term{loopy belief propagation}), but the induction fails at the
first cycle and the beliefs are approximations. The paper's claim is that its
inference is \emph{tree-exact}: the posterior graph is a tree, so no such approximation
is made.

\section{The chain posterior of the paper}

\subsection{Conditioning adds only unary factors}

The observation model is the coordinatewise channel of the paper's equation
\emph{channel}: given the clean chain $\va$, the coordinates of
$\vx=e^{-t}\va+\sqrt{\Dt}\vz$ are independent, each with density
$\normal{e^{-t}a_i}{\Dt}$ in $x_i$. By Bayes' rule,
\begin{equation}
  p_t(\va\mid\vx)=\frac{\Pz(\va)\prod_{i=1}^{\nsites}\like(x_i\mid a_i)}{\Pt(\vx)},\qquad
  \like(x_i\mid a_i)=\frac{1}{\sqrt{2\pi\Dt}}\exp\!\Bigl[-\frac{(x_i-e^{-t}a_i)^2}{2\Dt}\Bigr].
  \label{eq:bp-posterior}
\end{equation}
Compared with the prior, the numerator has one extra factor per site, each of scope
$\{i\}$: a \term{unary likelihood factor}. No factor couples two latent variables that
were not already coupled. The factor graph of the posterior is therefore the prior's
path with a unary factor hanging from every variable node; it is still a tree
(\cref{fig:bp-chain}). This is the whole content of the sentence ``conditioning leaves
the posterior factor graph a tree''. It would fail for a channel that mixed coordinates
(a non-diagonal noise covariance would add pairwise factors among the $a_i$, closing
cycles), and it holds for any prior whose graph is a tree, not only for a path.

\begin{figure}[t]
  \centering
  \begin{tikzpicture}[
    x=1cm, y=1cm, font=\small,
    lat/.style={circle, draw=bpInk, line width=0.5pt, fill=white, minimum size=6.2mm, inner sep=0pt},
    obs/.style={circle, draw=bpInk, line width=0.5pt, fill=bpFill, minimum size=6.2mm, inner sep=0pt},
    fac/.style={rectangle, draw=bpInk, line width=0.5pt, fill=bpInk, minimum size=2.2mm, inner sep=0pt},
    ufac/.style={rectangle, draw=bpInk, line width=0.5pt, fill=white, minimum size=2.2mm, inner sep=0pt},
    edge/.style={draw=bpInk, line width=0.5pt},
    gap/.style={inner sep=0pt, minimum size=6.2mm, text=bpInk},
    msg/.style={-{Latex[length=1.8mm, width=1.4mm]}, line width=0.8pt, draw=bpBlue},
    rowlab/.style={anchor=east, inner sep=0pt},
  ]
    \def\dx{1.55}
    \def\ya{1.5}\def\yl{0.75}\def\yx{0.0}
    \foreach \i/\al in {1/a_1, 2/a_2, 4/a_i, 6/a_\nsites} {
      \node[lat] (a\i) at (\dx*\i, \ya) {$\al$};
    }
    \foreach \i in {3,5} { \node[gap] (a\i) at (\dx*\i, \ya) {$\cdots$}; }
    \foreach \i/\j in {1/2,2/3,3/4,4/5,5/6} {
      \node[fac] (k\i) at ($(a\i)!0.5!(a\j)$) {};
      \draw[edge] (a\i) -- (k\i) -- (a\j);
    }
    \foreach \i/\xl in {1/x_1, 2/x_2, 4/x_i, 6/x_\nsites} {
      \node[ufac] (l\i) at (\dx*\i, \yl) {};
      \node[obs]  (x\i) at (\dx*\i, \yx) {$\xl$};
      \draw[edge] (a\i) -- (l\i) -- (x\i);
    }
    \foreach \i in {3,5} { \node[gap] (x\i) at (\dx*\i, \yx) {$\cdots$}; }
    \node[fac] (m1) at (\dx*1-0.75, \ya) {};
    \draw[edge] (m1) -- (a1);
    \node[anchor=east, inner sep=1pt, font=\footnotesize] at (\dx*1-0.85, \ya) {$\mu$};
    \node[rowlab] at (\dx-1.35, \ya+0.55) {$\kernel(a_i\!\mid\! a_{i-1})$};
    \node[rowlab] at (\dx-1.35, \yl) {$\like(x_i\!\mid\! a_i)$};
    \node[rowlab] at (\dx-1.35, \yx) {$\vx_t$};
    \draw[msg] ($(a1)+(0.35,0.55)$) -- ($(a4)+(-0.4,0.55)$)
      node[midway, above=1pt, text=bpBlue, inner sep=1pt] {$\lmsg$};
    \draw[msg] ($(a6)+(-0.35,0.55)$) -- ($(a4)+(0.4,0.55)$)
      node[midway, above=1pt, text=bpBlue, inner sep=1pt] {$\rmsg$};
  \end{tikzpicture}
  \caption{The factor graph of the chain posterior, drawn as in the paper's Figure 1.
  Latent sites (white circles) are joined by the transition factors $\kernel$ (filled
  squares) and, at the left end, by the initial factor $\mu$. Conditioning attaches one
  unary likelihood factor (open square) per site, joining $a_i$ to its own observation
  $x_i$ (shaded circle), and adds no edge among the latent variables. The graph is a
  tree, so the forward sweep $\protect\lmsg$ and the backward sweep $\protect\rmsg$ of the sum-product
  algorithm are exact; they meet at the query site $i$.}
  \label{fig:bp-chain}
\end{figure}

\subsection{From the general rules to the two sweeps}

Specialise \cref{eq:bp-v2f,eq:bp-f2v} to \cref{fig:bp-chain}. Each latent variable $a_i$
has at most three neighbouring factors: the transition factor on its left, the
transition factor on its right, and its unary likelihood. Write $\kernel_{i}$ for the
factor between $a_{i-1}$ and $a_i$ (with $\kernel_1=\mu$ the unary initial factor).
Going left to right, the message that $\kernel_{i}$ sends into $a_i$ is, by
\cref{eq:bp-f2v}, an integral over $a_{i-1}$ of $\kernel(a_i\mid a_{i-1})$ times the
message $a_{i-1}$ sends into $\kernel_i$; and that message is, by \cref{eq:bp-v2f}, the
product of what $a_{i-1}$ receives from its other two neighbours, its own unary factor
$\like(x_{i-1}\mid\cdot)$ and the transition factor on its left. Naming
$\lmsg_i(a)=m_{\kernel_i\to a_i}(a)$, this reads
\begin{equation}
  \lmsg_1(a)=\mu(a),\qquad
  \lmsg_i(a)=\int\kernel(a\mid a')\,\like(x_{i-1}\mid a')\,\lmsg_{i-1}(a')\dd a',\quad i\ge2.
  \label{eq:bp-forward}
\end{equation}
Going right to left, name $\rmsg_i(a)=m_{\kernel_{i+1}\to a_i}(a)$ the message that the
transition factor on the right of $a_i$ sends into $a_i$; for the last site there is
no such factor and $\rmsg_\nsites\equiv1$ (an empty product). The same two rules give
\begin{equation}
  \rmsg_\nsites(a)=1,\qquad
  \rmsg_i(a)=\int\kernel(a''\mid a)\,\like(x_{i+1}\mid a'')\,\rmsg_{i+1}(a'')\dd a'',\quad i\le\nsites-1.
  \label{eq:bp-backward}
\end{equation}
These are, up to the proportionality signs discussed in \cref{sec:bp-normalisation},
the two lines of the paper's equation \emph{bp}. The paper's convention for
$\lmsg_1$ differs from \cref{eq:bp-forward} in one cosmetic respect: it writes
$\lmsg_1\equiv1$ and treats $\mu$ as if it were part of the first site's unary factor.
The code follows \cref{eq:bp-forward}, with \code{L[0] = exp(log\_mu)} and
\code{R[-1] = 1} in \code{src/markov\_score/bp\_grid.py}. Note the asymmetry of the
kernel's arguments: forward, the integration variable is the \emph{conditioning}
argument of $\kernel$; backward, it is the \emph{conditioned} one. The transition
matrix on a grid, $\kernel[k,l]=\kernel(\gridpt{k}\mid\gridpt{l})$ with column index
the parent, enters the forward sweep as $\kernel$ and the backward sweep as
$\kernel^{\!\top}$, which is exactly what the code does.

Two conventions are worth fixing, because both appear in the literature. The messages
\cref{eq:bp-forward,eq:bp-backward} \emph{exclude} the unary factor of the site they
arrive at; they are the factor-to-variable messages of the general algorithm. The
alternative convention (the ``filtered'' forward variable of the HMM literature)
multiplies the local likelihood in and would be $\lmsg_i\like$. The paper and the code
use the excluded form, so that the local factor appears once, explicitly, in every
formula below.

\subsection{Site marginals, pairwise marginals and the evidence}

By \cref{eq:bp-beliefs}, the belief at site $i$ is the product of what it receives from
its three neighbours,
\begin{equation}
  p_t(a_i\mid\vx)=\frac{\lmsg_i(a_i)\,\like(x_i\mid a_i)\,\rmsg_i(a_i)}{Z_i},\qquad
  Z_i=\int\lmsg_i\like(x_i\mid\cdot)\rmsg_i\dd a_i,
  \label{eq:bp-site}
\end{equation}
and by \cref{thm:bp-exact} it is the exact posterior marginal. The paper's phrase
``site marginals $\propto\lmsg_i\like\rmsg_i$'' is this line. The belief at the
transition factor between $a_{i-1}$ and $a_i$ is
\begin{equation}
  p_t(a_{i-1},a_i\mid\vx)=\frac{\lmsg_{i-1}(a_{i-1})\,\like(x_{i-1}\mid a_{i-1})\,
    \kernel(a_i\mid a_{i-1})\,\like(x_i\mid a_i)\,\rmsg_i(a_i)}{Z_{i-1,i}},
  \label{eq:bp-pair}
\end{equation}
the exact \term{pairwise marginal}, whose two single-site margins are \cref{eq:bp-site}
at $i-1$ and at $i$ (local consistency). The pairwise marginals are what the
expectation--maximisation chapter (\cref{ch:em-through-noise}) sums into the expected
transition statistic $\Xis$; they come from the same two sweeps at no extra cost beyond
the outer products.

The normaliser is available at any site, $Z_i=\Pt(\vx)$ for every $i$, because every
$Z_i$ equals $\int\Pz(\va)\prod_j\like(x_j\mid a_j)\dd\va$ by the proof of
\cref{thm:bp-exact}; the same holds for every $Z_{i-1,i}$. In particular the forward
sweep alone gives it, as $Z_\nsites=\int\lmsg_\nsites\like(x_\nsites\mid\cdot)\dd a$ once
$\rmsg_\nsites\equiv1$. This number, $\log\Pt(\vx)$, is the \term{log-evidence}: the
marginal log-likelihood of one noised sequence under the chain prior, which the paper's
Section 4 maximises over the parameters of the kernel and which its Appendix E uses,
summed over a held-out bundle, as an oracle-free model-selection criterion.

\subsection{From the marginal to the score}

The paper needs one number per site, the posterior mean
$\pmeansite{i}=\int a_i\,p_t(a_i\mid\vx)\dd a_i$, because Tweedie's identity
(the chapter on the Ornstein--Uhlenbeck channel) converts it into the score,
$\score(\vx,t)_i=-(x_i-e^{-t}\pmeansite{i})/\Dt$. Thus the whole of ``the exact score''
is: two sweeps, $\nsites$ one-dimensional integrals for the means, and one subtraction
per site. The word ``exact'' refers to \cref{thm:bp-exact}; what is not exact is the
numerical representation of the messages, which is the subject of
\cref{ch:grid-quadrature}.

\subsection{Normalisation and the log-normaliser}
\label{sec:bp-normalisation}

Messages are defined up to a positive constant: multiplying $\lmsg_i$ by $c$ multiplies
every later forward message and every belief involving it by $c$, and the
normalised beliefs \cref{eq:bp-site,eq:bp-pair} are unchanged. This is why the paper
writes $\propto$ in its recursion. Numerically, unnormalised messages on a long chain
underflow or overflow (each step multiplies by a likelihood of order
$1/\sqrt{2\pi\Dt}$ and by a kernel of order $1/\sqrt{\ivar}$), so an implementation
rescales each message to unit mass and keeps the discarded constants:
\begin{equation}
  \tilde{\lmsg}_{i}=\frac{\lmsg_i}{c_i},\qquad c_i=\int\lmsg_i\dd a,\qquad
  \log\Pt(\vx)=\sum_{i=2}^{\nsites}\log c_i+\log\int\tilde{\lmsg}_\nsites\like(x_\nsites\mid\cdot)\dd a,
  \label{eq:bp-logZ}
\end{equation}
where each $c_i$ is computed from the \emph{already rescaled} previous message, so
that the product of the $c_i$ is the total constant removed. The log-evidence is thus
assembled from a sum of logarithms and never as a product of small numbers.
\Cref{ch:grid-quadrature} gives the grid version, including the constants of the
likelihood factors that the implementation removes for stability and restores at the
end.

\subsection{Cost}

On a grid of $\ngrid$ points, a message is a vector of length $\ngrid$ and a
transition integral is a matrix--vector product with the $\ngrid\times\ngrid$ matrix
$\kernel[k,l]$, costing $O(\ngrid^2)$. The forward sweep performs $\nsites-1$ of them,
the backward sweep the same, and the beliefs cost $O(\nsites\ngrid)$. The total is
$O(\nsites\ngrid^2)$ per sequence and per noise level, as the paper states in its
conclusion, and it is linear in the chain length: this is why the length-extrapolation
study of the paper's Section 5 costs nothing extra at $\nsites=256$. For a batch of
$B$ sequences the matrix--vector products become matrix--matrix products
$(\ngrid\times\ngrid)(\ngrid\times B)$, which is how \code{grid\_bp\_batch} in
\code{src/markov\_score/bp\_grid.py} organises the computation.

\section{Three names for one recursion}

\paragraph{Forward--backward for hidden Markov models.}
When the state space is a finite set of $\ngrid$ symbols, \cref{eq:bp-forward,eq:bp-backward}
are the forward and backward recursions of Baum and colleagues (1970) and Rabiner's
tutorial (1989), usually written for the filtered variable
$\alpha_i(k)=P(a_i=k,x_1,\dots,x_i)$ and the backward variable
$\beta_i(k)=P(x_{i+1},\dots,x_\nsites\mid a_i=k)$. In our notation
$\alpha_i=\lmsg_i\like(x_i\mid\cdot)$ and $\beta_i=\rmsg_i$; the posterior at a site is
$\alpha_i\beta_i/\Pt(\vx)$ and the pairwise posterior is
$\alpha_{i-1}\kernel\like\beta_i/\Pt(\vx)$, both as above. The pairwise posteriors summed
over sites are the ``expected transition counts'' of the Baum--Welch algorithm, and
their continuous analogue on the grid is exactly the paper's $\Xis$. The paper's
recursion is therefore not a new algorithm; the paper says so in its related work,
citing Baum and Rabiner.

\paragraph{The Kalman smoother.}
When $\mu$, $\kernel$ and $\like$ are all Gaussian, every message is Gaussian (a
product of Gaussians is Gaussian, and the convolution of a Gaussian with a Gaussian
kernel is Gaussian), so each message is carried by two numbers, a mean and a variance,
or equivalently a precision and an information vector. The forward sweep is then the
Kalman filter and the combination of forward and backward messages is the
Rauch--Tung--Striebel smoother. \Cref{ch:gaussian-closure-lmmse} writes out the two
moment maps and proves the paper's Proposition 1, which says that projecting every
message onto a Gaussian, for \emph{any} innovation law with the same variance,
reproduces this smoother. The implementation of the Gaussian case is
\code{gaussian\_chain\_bp} in \code{src/markov\_score/bp\_gaussian.py}.

\paragraph{Functional message passing and Kitagawa's grid smoother.}
When the state is continuous and non-Gaussian, the messages are genuine functions on
$\Reals$ and have no finite parametrisation. The paper calls the recursion
``functional message passing'' to stress that the algorithm is defined at the level of
functions, before any discretisation, and ``tree-exact inference'' to stress that the
only approximation is in representing those functions. Representing them by their
values on a grid and replacing integrals by quadrature is Kitagawa's non-Gaussian
smoother (1987), which the paper also cites. The phrase ``the score'' throughout the
paper means: the posterior mean obtained by this recursion, evaluated on a grid whose
adequacy is measured (\cref{ch:grid-quadrature}), inserted into Tweedie's identity.

\section{A worked example: three sites, three states}
\label{sec:bp-example}

Take $\nsites=3$ sites and a three-point alphabet $\{\gridpt{1},\gridpt{2},\gridpt{3}\}=\{-1,0,1\}$
with no quadrature weights (a finite-state model, so that every integral is a sum).
Let the initial law, the transition matrix and the three likelihood rows be
\[
  \mu=\begin{pmatrix}0.25\\0.5\\0.25\end{pmatrix},\qquad
  \kernel=\begin{pmatrix}0.6&0.2&0.1\\0.3&0.6&0.3\\0.1&0.2&0.6\end{pmatrix},\qquad
  \ell_1=\begin{pmatrix}0.2\\0.5\\0.3\end{pmatrix},\;
  \ell_2=\begin{pmatrix}0.1\\0.3\\0.6\end{pmatrix},\;
  \ell_3=\begin{pmatrix}0.4\\0.4\\0.2\end{pmatrix},
\]
where $\kernel[k,l]$ is the probability of moving from state $l$ (column) to state $k$
(row), so the columns sum to one, and $\ell_i[k]$ (the unary factor $\like(x_i\mid\gridpt{k})$ of the chapter, written as a plain vector here) is the likelihood of the $i$-th
observation given $a_i=\gridpt{k}$ (any positive numbers; they are not normalised over
$k$ and need not be).

\emph{Forward sweep.} $\lmsg_1=\mu$. The message into site 2 is
$\lmsg_2=\kernel(\lmsg_1\circ\ell_1)$, with $\circ$ the elementwise product:
$\lmsg_1\circ\ell_1=(0.05,\,0.25,\,0.075)$, and
\[
  \lmsg_2=\begin{pmatrix}0.6\cdot0.05+0.2\cdot0.25+0.1\cdot0.075\\
                         0.3\cdot0.05+0.6\cdot0.25+0.3\cdot0.075\\
                         0.1\cdot0.05+0.2\cdot0.25+0.6\cdot0.075\end{pmatrix}
        =\begin{pmatrix}0.0875\\0.1875\\0.1\end{pmatrix}.
\]
Then $\lmsg_2\circ\ell_2=(0.00875,\,0.05625,\,0.06)$ and
$\lmsg_3=\kernel(\lmsg_2\circ\ell_2)=(0.0225,\,0.054375,\,0.048125)$.

\emph{Backward sweep.} $\rmsg_3=(1,1,1)$. Then
$\rmsg_2=\kernel^{\!\top}(\rmsg_3\circ\ell_3)=\kernel^{\!\top}(0.4,0.4,0.2)$: the first
entry is $0.6\cdot0.4+0.3\cdot0.4+0.1\cdot0.2=0.38$, and altogether
$\rmsg_2=(0.38,\,0.36,\,0.28)$. Next $\rmsg_2\circ\ell_2=(0.038,\,0.108,\,0.168)$ and
$\rmsg_1=\kernel^{\!\top}(\rmsg_2\circ\ell_2)=(0.072,\,0.106,\,0.137)$.

\emph{Evidence.} From the end of the forward sweep,
$Z=\sum_k\lmsg_3[k]\ell_3[k]=0.0225\cdot0.4+0.054375\cdot0.4+0.048125\cdot0.2=0.040375$.
The same number is obtained at site 1, $\sum_k\mu[k]\ell_1[k]\rmsg_1[k]
=0.25\cdot0.2\cdot0.072+0.5\cdot0.5\cdot0.106+0.25\cdot0.3\cdot0.137=0.040375$, and at
site 2, as \cref{thm:bp-exact} requires.

\emph{Site marginals.} At site 1 the unnormalised belief is
$\mu\circ\ell_1\circ\rmsg_1=(0.0036,\,0.0265,\,0.010275)$; dividing by $Z$ gives
$p(a_1\mid\vx)=(0.0892,\,0.6564,\,0.2545)$ and posterior mean
$-0.0892+0.2545=0.1653$. At site 2 the unnormalised belief is
\[
  \lmsg_2\circ\ell_2\circ\rmsg_2=(0.003325,\,0.02025,\,0.0168),
\]
so that the normalised marginal is $p(a_2\mid\vx)=(0.0824,\,0.5015,\,0.4161)$, with
mean $0.3337$. At site 3 the same steps give $p(a_3\mid\vx)=(0.2229,\,0.5387,\,0.2384)$
and the mean $0.0155$. Enumerating all $27$
configurations of $(a_1,a_2,a_3)$ and summing gives the same $Z$ and the same three
marginals, which is the finite version of \cref{thm:bp-exact}.

\emph{Pairwise marginal.} On the edge $(1,2)$, \cref{eq:bp-pair} gives the matrix with
entries
\[
  p(a_2=\gridpt{k},a_1=\gridpt{l}\mid\vx)=\rmsg_2[k]\,\ell_2[k]\,\kernel[k,l]\,\lmsg_1[l]\,\ell_1[l]/Z;
\]
for instance the entry $(k,l)=(2,2)$, staying at $0$, is
$0.108\cdot0.6\cdot0.25/0.040375=0.4012$, and the full matrix is
\[
  p(a_2=\gridpt{k},a_1=\gridpt{l}\mid\vx)=
  \begin{pmatrix}0.0282&0.0471&0.0071\\0.0401&0.4012&0.0602\\0.0208&0.2081&0.1872\end{pmatrix},
\]
whose column sums are the site-1 marginal and whose row sums are the site-2 marginal,
to the rounding shown. This is the object that \cref{ch:em-through-noise} accumulates
into $\Xis$: summing the two edge matrices of this chain gives a $3\times3$ matrix of
total mass $2$, the number of transitions.

\where{The paper, Section 2 (``the score is a posterior mean, and the posterior is a
tree'', Figure 1) and Appendix B (the recursion \emph{bp}); Section 4 (the pairwise
marginals behind $\Xis$); Appendix E (the log-evidence as a selection criterion, E20).
Code: \code{src/markov\_score/bp\_grid.py} (\code{grid\_bp} and \code{grid\_bp\_batch}),
\code{src/markov\_score/em.py} (\code{e\_step}, which adds the pairwise accumulation to
the same sweeps), \code{src/markov\_score/bp\_gaussian.py} (the Gaussian special case),
and \code{paper/figures/fig\_mechanism.tex} (the figure mirrored above). Tests: the
brute-force enumeration test in \code{tests/test\_em\_bp.py};
\code{tests/test\_fisher\_bruteforce.py}, which checks the evidence and the pairwise
marginals against enumeration on a five-state, four-site chain; and
\code{tests/test\_gaussian\_bp\_equivalence.py}. Experiments: every registered
experiment, E01 to E21, scores against this recursion.}

\paragraph{References.}
J.~Pearl (1988), \emph{Probabilistic Reasoning in Intelligent Systems}, Morgan Kaufmann.
F.~R.~Kschischang, B.~J.~Frey and H.-A.~Loeliger (2001), ``Factor graphs and the sum-product
algorithm'', \emph{IEEE Transactions on Information Theory} 47(2).
L.~E.~Baum, T.~Petrie, G.~Soules and N.~Weiss (1970), ``A maximization technique occurring
in the statistical analysis of probabilistic functions of Markov chains'', \emph{Annals of
Mathematical Statistics} 41(1).
L.~R.~Rabiner (1989), ``A tutorial on hidden Markov models and selected applications in
speech recognition'', \emph{Proceedings of the IEEE} 77(2).
G.~Kitagawa (1987), ``Non-Gaussian state-space modeling of nonstationary time series'',
\emph{Journal of the American Statistical Association} 82(400).
M.~M\'ezard and A.~Montanari (2009), \emph{Information, Physics, and Computation}, Oxford
University Press, Chapter 14.
